\chapter{Apr.~7 --- Classification of Curves}

\section{Higher Genus Curves}

\begin{remark}
  Recall that a smooth plane curve
  $C_d \subseteq \PP^2$ of
  degree $d$ has genus
  $(d - 1)(d - 2) / 2$. This does not cover
  most $g \ge 2$. One may also try to
  consider complete intersections, but again
  this does not cover everything.
  The full answer is the following:
  \begin{enumerate}
    \item For $g \ge 2$, there exists an
      irreducible quasiprojective variety
      $\mathcal{M}_g$
      of dimension $3g - 3$ parametrizing
      all smooth projective genus $g$ curves.
    \item If $C \in \mathcal{M}_g$,
      then either
      \begin{itemize}
        \item there exists a 2-to-1
          cover $C \to \PP^1$
          (\emph{hyperelliptic curves}), or
        \item $\OO_C(K_C)$
          is very ample and induces an
          embedding $C \hookrightarrow \PP^{g - 1}$
          (\emph{canonically embedded curves}).
      \end{itemize}
  \end{enumerate}
  To show (2) (we will not prove (1)), we need
  to understand \emph{finite morphisms}.
\end{remark}

\section{Finite Morphisms}

\begin{definition}
  A morphism of varieties $f : X \to Y$ is
  \emph{finite} if there exists an open
  affine cover $\{U_i\}$ of $Y$ such that
  $V_i = f^{-1}(U_i)$ is affine and
  $\OO_Y(U_i) \to \OO_X(V_i)$ is a finite
  morphism of rings (i.e.
  $\OO_X(V_i)$ is a finite $\OO_Y(U_i)$-module).
\end{definition}

\begin{example}
  Let $d > 0$ and consider the morphism
  \begin{align*}
    \Affine^1_x &\longrightarrow \Affine^1_y \\
    x &\longmapsto x^d.
  \end{align*}
  This morphism is finite, as
  $k[y] \to k[x]$ is given by $y \mapsto x^d$,
  which is finite
  (note that
  $k[x]$ is generated as a $k[y]$-module by
  $1, x, \dots, x^{d - 1}$).
\end{example}

\begin{remark}
  We have the following properties of
  finite morphisms:
  \begin{enumerate}
    \item If $f : X \to Y$ is finite and
      $U \subseteq Y$ is open affine, then
      $V = f^{-1}(U)$ is open affine and
      $\OO_Y(V) \to \OO_X(U)$ is finite.
      (See Mustata 5.3.2 for a proof.)
    \item A closed embedding
      $Z \hookrightarrow Y$ is finite.

      To see this, note that
      it suffices to prove this on an
      open cover, so we may assume
      $Y$ is affine. Set $A = \OO_Y(Y)$.
      Then $Z = V(I_Z)$ with $I_Z \subseteq A$
      and $A = \OO_Y(Y) \to \OO_Z(Z) = A / I_Z$,
      which is finite.
    \item If $X \to Y$ is finite, then
      $\# f^{-1}(\{y\}) < \infty$ for any
      $y \in Y$.

      To see this, note that this is local in
      a neighborhood of $y$, so we may assume
      that $Y$ is affine. Then by (1),
      $X$ is also affine. Then
      $B := \OO_Y(Y) \to \OO_X(X) =: C$
      is finite. Now
      \[
        A(f^{-1}(\{y\}))
        = C / I_{f^{-1}(\{y\})}
        = (C / \m_y C)^{\mathrm{red}}
        = (C \otimes_B B / \m_y B)^{\mathrm{red}}
      \]
      where $R^{\mathrm{red}}$ denotes
      the reduced ring corresponding to $R$.
      Now $C \otimes_B B / \m_y B$ is a
      finite $B / \m_y B$-module and
      $B / \m_y \cong k$, so we see that
      $\dim_k A(f^{-1}(\{y\})) < \infty$.
      Thus it suffices to show:
      \begin{quote}
        \textbf{Claim.} If $X$ is an affine
        variety and $\dim_k A(X) < \infty$,
        then $X$ is a finite set of points.

        \begin{proof}[Proof of claim]
          We may assume that $X$ is
          irreducible.\footnote{At this point, this is essentially the statement that a finite integral domain is a field, but with a finite-dimensional vector space instead. Either way, the key property is that an injective (linear) map is automatically also surjective.} Since $\dim_k A(X) < \infty$
          and $A(X)$ is a domain, for any
          $0 \ne f \in A(X)$, we have an
          injective linear map given by
          \begin{align*}
            A(X) &\longrightarrow A(X) \\
            1 &\longmapsto f.
          \end{align*}
          By finite dimensionality, this map
          is also surjective, so $f$ has an inverse.
          Thus $A(X)$ is a field, and hence
          $X$ is a single point.
        \end{proof}
      \end{quote}
  \end{enumerate}
\end{remark}

\begin{remark}
  Morphisms with finite fibers are often not
  finite. For example, the inclusion
  $\Affine^1 \setminus \{0\} \hookrightarrow \Affine^1$
  has finite fibers, but
  we have $k[x] = A(\Affine^1) \to A(\Affine^1 \setminus \{0\}) = k[x^{\pm 1}]$.
\end{remark}

\begin{theorem}\label{thm:finite-morphism-finite-fibers}
  A morphism of complete varieties
  $f : X \to Y$ with finite fibers is finite.
\end{theorem}

\begin{proof}
  See Mustata 14.1.8.
\end{proof}

\section{Finite Morphisms of Curves}

\begin{remark}
  Assume $f : X \to Y$ is a surjective
  morphism of smooth projective curves. Note
  that such a morphism always has
  finite fibers (as $\dim X = \dim Y = 1$),
  so $f$ is finite by
  Theorem \ref{thm:finite-morphism-finite-fibers}.
\end{remark}

\begin{definition}
  Define the \emph{degree} of $f$ to be
  $\deg f = \deg (K(X) / K(Y))$,
  where $f^* : K(Y) \hookrightarrow K(X)$.
  Note that since
  $\trdeg_k K(Y) = \trdeg_k K(X) = 1$,
  this extension is finite.
\end{definition}

\begin{definition}
  We say that $f$ is \emph{separable}
  if $K(X) / K(Y)$ is separable (recall that this
  condition always
  holds when $\Char k = 0$.)
\end{definition}

\begin{example}
  The morphism
    $\Affine^1_{\overline{\F}_p} \to \Affine^1_{\overline{\F}_p}$
  given by $x \to x^p$
  is non-separable.
\end{example}

\begin{prop}
  If $D$ is a divisor on $Y$, then
  $\deg f^* D = (\deg f)(\deg D)$. Here,
  $f^* D$ is defined by pulling back
  local equations (recall that every divisor
  is Cartier when $Y$ is smooth).
\end{prop}

\begin{proof}
  As the formula is linear in $D$, we may
  assume that $D = P$ for some $P \in Y$.
  Fix an open affine $P \in U \subseteq Y$, and
  set $V = f^{-1}(U)$. Now consider
  \[
    A = \OO_Y(U)
    \overset{f^*}{\longrightarrow}
    \OO_X(V) = B,
  \]
  which is finite by our assumptions.
  Let $\m_P \subseteq A$ be the ideal
  of functions vanishing at $P$. Now
  $B_{\m_P}$ is a finite $A_{\m_P}$-module,
  $A_{\m_P}$ is a DVR, and $B_{\m_P}$ is
  torsion-free, so $B_{\m_P}$ is a finite
  rank free $A_{\m_P}$-module.
  So we have $B_{\m_P} \cong A_{\m_P}^{\oplus r}$ for some
  $r \ge 0$, and $r = d := \deg f$ by localizing
  away from $0$. So
  \pagebreak
  \[
    d = \dim_k (B_{\m_P} / \m_P B_{\m_P})
    = \dim_k H^0(X, \OO_{f^* P})
    = \chi(X, \OO_{f^* P}).
  \]
  Now we can consider the short exact sequence
  \[
    \begin{tikzcd}
      0 \ar[r] & \OO_X(-f^* P) \ar[r] & \OO_X \ar[r] & \OO_{f^* P} \ar[r] & 0
    \end{tikzcd}
  \]
  from which we see that
  $\chi(X, \OO_{f^* P}) = \chi(X, \OO_X) - \chi(X, \OO_X(-f^* P)) = \deg f^* P$
  by Riemann-Roch. So
  \[
    \deg f^* P = d = (\deg f)(\deg P)
  \]
  since $\deg P = 1$,
  which is the desired formula.
\end{proof}

\begin{example}
  Let $f : \Affine^1 \to \Affine^1$
  be given by $x \mapsto x^2$. Then
  \[
    f^*([0] + [1])
    = 2[0] + [1] + [-1].
  \]
\end{example}

\begin{remark}
  How do we compare
  $\genus(X)$ and $\genus(Y)$? We have
  an exact sequence
  \[
    \begin{tikzcd}
      f^*\Omega_Y
      \ar[r, "\alpha"] & \Omega_X
      \ar[r] & \Omega_{X/Y} \ar[r] & 0
    \end{tikzcd} \tag{$*$}
  \]
  \begin{quote}
    \textbf{Claim.} If $f$ is separable,
    then $(*)$ is short exact.

    \begin{proof}
      Since $f^* \Omega_Y$ and $\Omega_X$
      are invertible, it suffices to show
      that $\alpha \ne 0$. Since
      $f$ is separable,
      we know that $\Omega_{K(X) / K(Y)} = 0$.
      Taking stalks at $p \in X$ and
      setting $q = f(p)$, we get
      \[
        \begin{tikzcd}
          \Omega_{\OO_{Y, q} / k}
          \otimes \OO_{X, p}
          \ar[r, "\alpha_p"] &
          \Omega_{\OO_{X, p}}
          \ar[r] & 
          \Omega_{\OO_{X, p} / \OO_{Y, q}}
          \ar[r] & 0
        \end{tikzcd}
      \]
      and applying $- \otimes K(X)$ gives
      \[
        \begin{tikzcd}[row sep=small]
          \Omega_{K(Y) / k} \otimes K(X)
          \ar[r, "\alpha_{K(X)}"] &
          \Omega_{K(X) / k}
          \ar[r] & \Omega_{K(X) / K(Y)} \ar[r] & 0 \\
                 & & 0 \ar[u, equal]
        \end{tikzcd}
      \]
      Thus $\alpha_{K(X)} \ne 0$, so
      $\alpha_P \ne 0$, and hence $\alpha \ne 0$.
    \end{proof}
  \end{quote}
  So now we assume that $f$ is separable
  and study $\Omega_{X / Y}$.
\end{remark}
